\section{Conformal expansion and the missing two powers}
\label{sec:jacobian}
%======================================================================

The square map explains the scale but not the sign.  Its Jacobian supplies the required two powers of dilution; the later Gram analysis must still explain cancellation.

Consider the continuous map
\begin{equation}
 \Phi(c,q)
 :=
 \left(
 x(c,q),y(c,q)
 \right)
 =
 \left(
 cq,\frac{q^2-c^2}{2}
 \right).
 \label{eq:Phi}
\end{equation}
Its differential is
\begin{equation}
 D\Phi(c,q)
 =
 \begin{pmatrix}
 q & c\\
 -c & q
 \end{pmatrix}.
 \label{eq:D-Phi}
\end{equation}
The two columns are orthogonal and have equal length $\sqrt{c^2+q^2}$.  Consequently,
\begin{equation}
 D\Phi(c,q)^{\mathsf T}D\Phi(c,q)
 =(c^2+q^2)I.
 \label{eq:conformal-metric}
\end{equation}
Thus $\Phi$ is conformal: locally it rotates and expands every direction by the same factor
\begin{equation}
 \lambda(c,q)=\sqrt{c^2+q^2}.
 \label{eq:length-scale}
\end{equation}
The area Jacobian is
\begin{equation}
 \boxed{
 J(c,q)=\det D\Phi(c,q)=c^2+q^2=2|w|.
 }
 \label{eq:Jacobian}
\end{equation}
By AM--GM,
\begin{equation}
 c^2+q^2\ge2cq=2x.
 \label{eq:Jacobian-lower}
\end{equation}
Therefore, in the square strip $x\asymp n^2$,
\begin{equation}
 \lambda(c,q)\gtrsim n,
 \qquad
 J(c,q)\gtrsim n^2.
 \label{eq:square-scale-expansion}
\end{equation}

This is the two-dimensional version of the vertical spacing estimate.  Unit-scale separation in the factor plane becomes scale-$n$ separation in the squared plane, while unit factor-lattice area becomes scale-$n^2$ area.  The coherent square-prefix ceiling loses exactly two powers if the relevant kernel neighborhoods have bounded squared-space area after localization.

\begin{caution}[What the Jacobian does not prove]
The M\"obius point cloud is discrete and arithmetically restricted, and the square-prefix kernel is not Euclidean convolution.  The Jacobian therefore does not by itself prove a signed variance estimate.  It proves the exact geometric dilution scale that a valid operator argument must inherit.
\end{caution}

\subsection{Two-dimensional area versus one-dimensional cofactor curves}
\label{subsec:curve-versus-area}

The conformal Jacobian governs two-dimensional regions in the continuous $(c,q)$ factor plane.  The arithmetic measure is not two-dimensional: for each fixed cofactor $c$, it is supported on the one-dimensional prime-sampled curve
\[
 q\longmapsto
 \Phi(c,q)
 =
 \left(cq,\frac{q^2-c^2}{2}\right),
 \qquad
 q\ \text{prime},\quad q>\Pplus(c).
\]
Along that curve,
\begin{equation}
 \left\lVert\frac{\partial\Phi}{\partial q}(c,q)\right\rVert
 =
 \sqrt{c^2+q^2}.
 \label{eq:curve-length-expansion}
\end{equation}
Thus one-dimensional arclength is also expanded, but the number and arithmetic placement of samples are governed by primes rather than by the area Jacobian.  A two-dimensional density estimate cannot be applied directly to a single cofactor curve.

The outer curve $c=1$ makes this obstruction explicit.  Its contribution to the $n$th complete prefix is
\begin{equation}
 S_n^{(1)}
 =
 -\pi(X_n),
 \label{eq:prime-curve-prefix}
\end{equation}
because every prime $q\le X_n$ has canonical pair $(1,q)$ and M\"obius weight $-1$.  The prime number theorem therefore predicts
\begin{equation}
 \sum_{n\le N}|S_n^{(1)}|^2
 \sim
 \frac{N^5}{20\log^2 N},
 \label{eq:prime-curve-asymptotic}
\end{equation}
up to the usual slowly varying correction.  The prime curve by itself is not on the cubic scale.  Cancellation between $c=1$ and the aggregate of the curves $c\ge2$ is indispensable.

This also shows why a bound of the form
\[
 \sum_{q\le Q}e\!\left(\alpha\frac{q^2-c^2}{2}\right)
 \ll Q^{1/2+\eps}
\]
cannot hold uniformly in $\alpha$: at $\alpha=0$ the left side is $\pi(Q)+O(1)$.  The zero and major-arc modes must be cancelled across cofactor curves by the M\"obius weights $\mu(c)$ and the exact smooth--transport architecture.  Oscillatory estimates are relevant only after that coherent component has been isolated.

\subsection{Small-modulus prime resonances}
\label{subsec:exact-prime-resonances}

The coherent rational modes are more rigid than a generic major-arc heuristic suggests.

\begin{theorem}[Prime squares modulo $24$]
\label{thm:prime-square-mod24}
If $q>3$ is prime, then
\begin{equation}
 \boxed{q^2\equiv1\pmod{24}.}
 \label{eq:prime-square-mod24}
\end{equation}
\end{theorem}

\begin{proof}
The prime $q$ is odd, so $q^2\equiv1\pmod8$.  Also $3\nmid q$, so $q\equiv\pm1\pmod3$ and $q^2\equiv1\pmod3$.  Since $8$ and $3$ are coprime, the Chinese remainder theorem gives \cref{eq:prime-square-mod24}.
\end{proof}

\begin{corollary}[Perfect coherence on the $2$--$3$ resonance lattice]
\label{cor:perfect-prime-coherence}
Let $m\in\mathbb Z$ and put $\alpha=m/12$.  For every prime $q>3$,
\begin{equation}
 e\!\left(\alpha\frac{q^2}{2}\right)
 =e\!\left(\frac{m q^2}{24}\right)
 =e\!\left(\frac m{24}\right).
 \label{eq:perfect-prime-phase}
\end{equation}
Therefore, if $Q>3$,
\begin{equation}
 \boxed{
 \sum_{Q<q\le2Q\atop q\ {\rm prime}}
 e\!\left(\alpha\frac{q^2}{2}\right)
 =e\!\left(\frac m{24}\right)
 \bigl(\pi(2Q)-\pi(Q)\bigr).
 }
 \label{eq:perfect-prime-sum}
\end{equation}
In particular, the magnitude equals the full prime count.
\end{corollary}

This exact coherence includes
\[
 \alpha\in\left\{0,\frac12,\frac13,\frac23,\frac14,\frac34,
 \frac16,\frac56\right\}.
\]
It is not a failure of the major/minor-arc philosophy.  It identifies a finite, explicit resonant component that must be reproduced by the major-arc model and then cancelled across cofactor curves.

\begin{proposition}[The rational quadratic phase has modulus $2r$]
\label{prop:phase-modulus-two-r}
Let $a,r,u\in\mathbb Z$ with $r\ne0$.  Then
\begin{equation}
 \frac{e\!\left(a(u+r)^2/(2r)\right)}
      {e\!\left(au^2/(2r)\right)}
 =(-1)^{ar}.
 \label{eq:phase-shift-r}
\end{equation}
Consequently $u\mapsto e(au^2/(2r))$ is always periodic modulo $2r$, but when $ar$ is odd it is not a well-defined function of $u\bmod r$.
\end{proposition}

\begin{proof}
Expanding the square gives
\[
 \frac{a(u+r)^2-au^2}{2r}=au+\frac{ar}{2}.
\]
The integer term $au$ disappears under $e(t)=e^{2\pi i t}$, while $e(ar/2)=(-1)^{ar}$.
\end{proof}

\begin{definition}[Corrected reduced quadratic prime factor]
\label{def:corrected-quadratic-factor}
For $r\ge1$ and $(a,r)=1$, define
\begin{align}
 G_{2r}^{\times}(a)
 &:={}
 \sum_{\substack{u\bmod 2r\\(u,2r)=1}}
 e\!\left(\frac{au^2}{2r}\right),
 \label{eq:corrected-reduced-Gauss}\\
 \gamma(a,r)
 &:={}
 \frac{G_{2r}^{\times}(a)}{\varphi(2r)}.
 \label{eq:normalized-corrected-factor}
\end{align}
\end{definition}

At $a=1,r=3$, the units modulo $6$ are $1$ and $5$, both with square $1$ modulo $6$.  Thus
\begin{equation}
 G_{6}^{\times}(1)=2e(1/6),
 \qquad
 \boxed{\gamma(1,3)=e(1/6),\quad |\gamma(1,3)|=1.}
 \label{eq:r3-corrected-factor}
\end{equation}
By contrast, the expression obtained by summing representatives modulo $3$ is
\[
 e(1/6)+e(4/6)=0.
\]
That vanishing is an artifact of using a modulus on which the phase is not well-defined.

%======================================================================
